% Proof inputs: C_F0082 and C_F0063 Appendix A for liberation contractions.
\section{Applications to liberation gradients}
\label{sec:liberation-gradients}

We apply Theorem~\ref{thm:intrinsic-analytic} to normal generators with $L^2$ liberation gradients and nonatomic spectral distributions.

We use the $L^2$ formulation of Voiculescu's liberation gradients in \cite[5.3--5.4 and Proposition~6.1]{Voi99}.
\begin{definition}\label{def:liberation-gradients}
Let $X=(X_1,\ldots,X_d)\in\M^d$ generate $\M$, and set
\[
 A_i:=W^*(X_i),\qquad B_i:=W^*(X_j:j\ne i),\qquad
 \mathcal P_i:=\operatorname{Alg}(A_i,B_i)
 \quad(1\leq i\leq d),
\]
where $B_1:=\C1$ if $d=1$.
We say that $X$ admits \emph{$L^2$ liberation gradients} if the derivations
\[
 \delta_i:\mathcal P_i\subset H\longrightarrow
          \mathcal P_i\otimes_{\C}\mathcal P_i,\qquad
 \delta_i(a):=a\otimes1-1\otimes a\quad(a\in A_i),\qquad
 \delta_i(b):=0\quad(b\in B_i),
\]
are well-defined and $1\otimes1\in\dom\delta_i^*$ for every $i$.
Here $\delta_i^*$ is the Hilbert space adjoint for the dense domain $\mathcal P_i\Omega$ and the target $H\bar\otimes H$.
The liberation gradients $j_i:=j(A_i:B_i):=-\delta_i^*(1\otimes1)\in H$ satisfy
\[
 j_i^*=-j_i,\qquad
 \tau(j_ip)=(\tau\otimes\tau)(\delta_ip)
 \quad(p\in\mathcal P_i).
\]
\end{definition}
This condition is equivalent to finiteness of the liberation Fisher information
$\varphi^*(A_1;\ldots;A_d):=\sum_i\|j_i\|_2^2$ in \cite[Definition~9.1]{Voi99}.

The following construction supplies the operators needed for Theorem~\ref{thm:intrinsic-analytic}.
\begin{lemma}\label{lem:gradient-liberation}
If $X$ admits $L^2$ liberation gradients in the sense of Definition~\ref{def:liberation-gradients}, then $X$ admits a normalized liberation dual system.
\end{lemma}
\begin{proof}
Fix $i$, and define on $\mathcal P_i$
\[
 \widehat\delta_i:=\mathrm i\delta_i,\qquad
 C_i^\circ(p\Omega):=(\id\otimes\tau)\widehat\delta_ip,\qquad
 \eta:=\mathrm i j_i.
\]
Then $\widehat\delta_i^*(1\otimes1)=\eta$.
The prescriptions on $A_i,B_i$ and the product rule show that $\widehat\delta_i$ is real:
$\widehat\delta_i(p^*)=(\widehat\delta_ip)^\dagger$, where $(u\otimes v)^\dagger:=v^*\otimes u^*$.
By \cite[Lemmas~4.6--4.7]{Mai15}, $\eta=\eta^*$, the algebraic tensor product $\mathcal P_i\otimes_{\C}\mathcal P_i$ is contained in $\dom\widehat\delta_i^*$, and
\[
 \langle C_i^\circ p,q\rangle=\langle p,q\eta-C_i^\circ q\rangle
 \qquad(p,q\in\mathcal P_i).
\]
In particular, both $\widehat\delta_i$ and $C_i^\circ$ are closable.
We adapt the argument in \cite[proof of Lemma~12]{Dab10}.
For $\widehat\delta_i p=\sum_r u_r\otimes v_r$, the prescriptions on $A_i,B_i$ and the product rule give
\[
 (\widehat\delta_i\otimes\id)\widehat\delta_i p-(\id\otimes\widehat\delta_i)\widehat\delta_i p
 =\mathrm i\sum_r u_r\otimes1\otimes v_r.
\]
Indeed, the mixed product-rule terms cancel, and the two sides agree on $A_i$ and $B_i$.
Since $C_i^\circ(\mathcal P_i)\subseteq\mathcal P_i$, applying $\id\otimes\tau\otimes\tau$ and using $\widehat\delta_i^*(1\otimes1)=\eta$ gives
\[
 (C_i^\circ)^2p=\sum_r u_r\tau(v_r\eta)+\mathrm iC_i^\circ p.
\]
The preceding integration by parts identity and the expansion of $\widehat\delta_i(p^*p)$ yield
\begin{equation}\label{eq:liberation-energy}
 \|C_i^\circ p\|_2^2
 =2\operatorname{Re}\langle C_i^\circ p,p\eta\rangle
  -\langle\widehat\delta_i(p^*p),1\otimes\eta\rangle
  -\langle\mathrm iC_i^\circ p,p\rangle.
\end{equation}

For $a=a^*\in\mathcal P_i$ with $\|a\|<1$, set $v:=a+\mathrm i(1-a^2)^{1/2}$.
The polynomial partial sums $v_n$ of its power series converge in operator norm.
The series for $\widehat\delta_i v_n$ and $\widehat\delta_i(v_n^*)$ converge in $H\bar\otimes H$, since $\|\widehat\delta_i(a^n)\|_2\leq n\|a\|^{n-1}\|\widehat\delta_i a\|_2$.
The product rule shows that $\widehat\delta_i(v_n^*v_n)$ converges, while $v_n^*v_n\to1$ in operator norm.
Closedness of $\overline{\widehat\delta_i}$ therefore gives $\widehat\delta_i(v_n^*v_n)\to0$.
Taking real parts and passing to the limit in \eqref{eq:liberation-energy}, with $\|v\|_2=1$, gives
\[
 \lim_n\|C_i^\circ v_n\|_2\leq2\|\eta\|_2+1.
\]
Apply the same argument to $v^*$; averaging via $a=(v+v^*)/2$, rescaling and splitting real and imaginary parts give, for some $K_i>0$,
\[
 \|C_i^\circ p\|_2\leq K_i\|p\|\qquad(p\in\mathcal P_i).
\]

The integration by parts identity shows that $\mathcal P_i\Omega\subseteq\dom(C_i^\circ)^*$ and
\[
 (C_i^\circ)^*(p\Omega)=p\eta-C_i^\circ(p\Omega),\qquad
 \|(C_i^\circ)^*(p\Omega)\|_2\leq(K_i+\|\eta\|_2)\|p\|
 \quad(p\in\mathcal P_i).
\]
Put $C_i:=\overline{C_i^\circ}$.
For $x\in\M$, Kaplansky density \cite[Theorem~2.3.1]{AP18} gives a uniformly bounded net $p_\lambda\in\mathcal P_i$ with $p_\lambda\Omega\to x\Omega$.
The estimate for $C_i^\circ p_\lambda$ gives a weakly convergent subnet in $H$.
The graph of $C_i$ is a closed linear subspace of $H\oplus H$ and hence is weakly closed, so $x\Omega\in\dom C_i$.
Applying the same argument to the bound for $(C_i^\circ)^*p_\lambda$ gives $x\Omega\in\dom C_i^*$.
Thus
\[
 \M\Omega\subseteq\dom C_i\cap\dom C_i^*.
\]
On $\mathcal P_i\Omega$, the product rule and the values of $\widehat\delta_i$ give
\[
 [C_i^\circ,X_j]=\mathrm i\delta_{ij}[X_i,P_\Omega].
\]
Since $\mathcal P_i\Omega$ is a core for $C_i$ and the right-hand side is bounded, approximation in the graph norm extends this identity to $\dom C_i$, with $X_j\dom C_i\subseteq\dom C_i$.
In particular it holds on $\M\Omega$.
Set $L_i:=\mathrm iC_i$.
Then both $\dom L_i$ and $\dom L_i^*$ contain $\M\Omega$, and
\[
 [X_j,L_i]|_{\M\Omega}=\delta_{ij}[X_i,P_\Omega]|_{\M\Omega}.
\]
This is a liberation dual system.
Lemma~\ref{lem:liberation-normalization} supplies a normalized system.
\end{proof}

For normal generators, nonatomicity supplies the remaining regularity hypothesis of Theorem~\ref{thm:intrinsic-analytic}.
\begin{proposition}\label{prop:liberation-gradients}
Let $X_1,\ldots,X_d\in\M$ be normal operators generating $\M$, each with nonatomic spectral distribution.
Suppose that $X$ admits $L^2$ liberation gradients in the sense of Definition~\ref{def:liberation-gradients}.
Then $\Delta(X)=d$, and $X$ admits a normalized liberation dual system.
For any such system, all the conclusions of Theorem~\ref{thm:intrinsic-analytic} hold.
In particular, for every $u\in L^0(\M)$ and every fraction pair $Q$ for $u$, we have $u\in\Div_{L^0(\M)}\A_X$ if and only if there exist $T_1,\ldots,T_d\in\F(H)$ such that
\[
 \sum_i[\rho(X_i),T_i]=\Gamma[Q,P_\Omega].
\]
For each fixed $Q$, the tuple $(T_1,\ldots,T_d)$ is unique whenever it exists.
\end{proposition}
\begin{proof}
Lemma~\ref{lem:gradient-liberation} gives a normalized liberation dual system.
Normality and nonatomicity imply that every $X_i-\lambda$, $\lambda\in\C$, is regular.
Theorem~\ref{thm:intrinsic-analytic} gives the criterion and uniqueness, and Proposition~\ref{prop:liberation-delta} gives $\Delta(X)=d$.
\end{proof}
